% R12 proofs of the whitened-innovation laws. Numerical checks: research/r7c/r12/verify_laws_output.txt.
\section{Proofs of the whitened-innovation laws}\label{sup:r12proofs}
Throughout, $w\sim\mathcal{CN}(0,I)$ denotes whitened clutter innovations (unit variance after dividing by $\sigma_e$), $g\sim\mathcal{CN}(0,1)$ the Swerling-1 amplitude, and $u$ the whitened target signature, so that the whitened episode is $w+gu\sim\mathcal{CN}(0,I+uu^{\mathsf H})$. $E$, $E_\pm$ are standard exponentials and $\Gamma_k$ a Gamma$(k,1)$ variable, all independent.

\subsection{Proposition (post-onset rank-one law)}
Write the detection event as $Q>0$ with $Q=x^{\mathsf H}Mx$, $x=w+gu$ and $M=\operatorname{diag}(-\beta I_m,1)$, $\beta=q^2/m$. As in Appendix~\ref{M-app:law} of the paper, $Q\overset{d}{=}\sum_k\lambda_kE_k$ with $\lambda_k$ the eigenvalues of $C^{1/2}MC^{1/2}$, $C=I+uu^{\mathsf H}$, which are those of $MC=M+(Mu)u^{\mathsf H}$. For a rank-one update of a diagonal matrix,
\[
\det(\lambda I-M-Muu^{\mathsf H})=\det(\lambda I-M)\Big(1-\sum_i\frac{M_{ii}|u_i|^2}{\lambda-M_{ii}}\Big)=(\lambda-1)(\lambda+\beta)^m\Big(1-\frac{a}{\lambda-1}+\frac{\beta b}{\lambda+\beta}\Big),
\]
with $a=|u_Y|^2$ and $b=\|u_H\|^2$. Hence $-\beta$ is an eigenvalue of multiplicity $m-1$ and the other two solve $(\lambda-1)(\lambda+\beta)-a(\lambda+\beta)+\beta b(\lambda-1)=0$, which is the quadratic of the proposition. Their product is $-\beta(1+a+b)<0$, so $\lambda_+>0>\lambda_-$. Then
\[
\PP\{Q>0\}=\PP\{\lambda_+E_+>|\lambda_-|E_-+\beta\Gamma_{m-1}\}=\mathbb E\,e^{-(|\lambda_-|E_-+\beta\Gamma_{m-1})/\lambda_+}=\frac{\lambda_+}{\lambda_++|\lambda_-|}\Big(1+\frac{\beta}{\lambda_+}\Big)^{-(m-1)}.
\]
For $b=0$ the roots are $1+a$ and $-\beta$, and the expression reduces to $(1+\beta/(1+a))^{-m}$. The signature of a persistent target follows from the whitened amplitudes in \eqref{M-eq:whitenedscr}: the onset innovation carries $S_w$, every later one $S_w|d|^2$.

\subsection{Corollary (visibility horizon)}
With $a=S_wD$ and $b=S_w(1+(\ell-1)D)$, $D=|d(\omega)|^2$, the linear coefficient is $B=1-\beta+S_w\kappa$, $\kappa=D-\beta\{1+(\ell-1)D\}$, and $\lambda_\pm=\{B\pm(B^2+4\beta(1+a+b))^{1/2}\}/2$ with $1+a+b=O(S_w)$. If $\kappa>0$, then $\lambda_+\sim S_w\kappa\to\infty$ while $\lambda_-$ stays bounded, so $P_{\rm d}\to1$. If $\kappa<0$, then $\lambda_+\to\beta(1+\ell D)/|\kappa|$ stays bounded while $\lambda_-\sim S_w\kappa\to-\infty$, so $\lambda_+/(\lambda_+-\lambda_-)\to0$ and $P_{\rm d}\to0$. The sign change $\kappa=0$ is $\ell=1+1/\beta-1/D$.

\subsection{Order-statistic scale}
Conditional on $G=|g|^2$, the history powers $|w_j+gu_j|^2$ are independent, exponential for clean innovations and non-central (with $2|w_j+gu_j|^2\sim\chi^2_2(2G|u_j|^2)$) for contaminated ones. The distribution function of their $k$-th smallest value $X_{(k)}$ is the Poisson-binomial probability that at least $k$ are below $x$, and $P_{\rm d}=\mathbb E_G\int\PP\{|w_Y+gu_Y|^2>q^2x\mid G\}\,dF_{X_{(k)}\mid G}(x)$. The code evaluates this by midpoint quadrature in the distribution function of $G$ and a dynamic program for the Poisson-binomial law. If $\ell\le m-k$ contaminated powers diverge, $X_{(k)}$ is the $k$-th smallest of the $m-\ell$ clean ones, which is independent of $w_Y$, and the OS-CFAR law with $m-\ell$ cells gives \eqref{M-eq:osstrong}. If $\ell>m-k$, $X_{(k)}$ diverges with the target while $|w_Y+gu_Y|^2/X_{(k)}$ tends to a constant below any practical threshold, as in Proposition~3.

\subsection{Dwell integration}
Under $H_0$ the event is $\Gamma_K>\tau\Gamma_m$, and conditioning on $\Gamma_m$ with the Poisson form of the Gamma tail gives $\sum_{i<K}\mathbb E[e^{-\tau\Gamma_m}(\tau\Gamma_m)^i/i!]=\sum_{i<K}\binom{m+i-1}{i}\tau^i(1+\tau)^{-(m+i)}$. With a Swerling-1 target the dwell form has one eigenvalue $c=1+b_D$ and $K-1$ unit ones. For $v>0$,
\[
\PP\{cE+\Gamma_{K-1}>v\}=\PP\{\Gamma_{K-1}>v\}+e^{-v/c}(1-1/c)^{-(K-1)}\PP\{\Gamma_{K-1}\le v(1-1/c)\}.
\]
Averaging over $v=\tau\Gamma_m$ uses $\mathbb E e^{-s\Gamma_m}=(1+s)^{-m}$ and $\mathbb E[e^{-s\Gamma_m}\PP\{\Gamma_{K-1}>\kappa\Gamma_m\mid\Gamma_m\}]=\sum_{i\le K-2}\binom{m+i-1}{i}\kappa^i(1+s+\kappa)^{-(m+i)}$ with $s+\kappa=\tau$, which gives \eqref{M-eq:dwellpd}.

\subsection{Self-normalized whitened Doppler}
With a known on-grid Doppler, the whitened dwell vector splits into its DFT coefficient on the target bin, $|\,\cdot\,|^2\sim(1+NS_w|d|^2)E$, and the orthogonal remainder $\Gamma_{N-1}$. The statistic $X/(X+\Gamma_{N-1})$ exceeds $t$ iff $X(1-t)>t\Gamma_{N-1}$, which has probability $\{1+t/((1-t)(1+NS_w|d|^2))\}^{-(N-1)}$; at $S_w=0$ this is $(1-t)^{N-1}$. For the maximum over the $N$ bins, let $I_j$ be independent exponentials with means $\mu_j$ and $J$ a set of bins with $|J|t<1$. Substituting $x_j=tT+y_j$ for $j\in J$, where $T=\sum_iI_i$, the map has Jacobian $1/(1-|J|t)$ and turns the exponent into a sum of independent exponentials, so
\[
\PP\{I_j>tT\ \forall j\in J\}=\frac{1}{1-|J|t}\prod_{i=1}^N\frac{1}{1+\theta_J\mu_i},\qquad\theta_J=\frac{t\sum_{j\in J}\mu_j^{-1}}{1-|J|t},
\]
and zero if $|J|t\ge1$. With all $\mu_i=1$ this is $(1-|J|t)^{N-1}$, and inclusion–exclusion gives Fisher's law $\sum_r(-1)^{r-1}\binom Nr(1-rt)^{N-1}$ (R.~A. Fisher, Proc.\ R.\ Soc.\ Lond.\ A 125, 1929). With one bin of mean $1+NS_w|d|^2$, inclusion–exclusion over sets that do or do not contain it gives $P_{\rm d}$ (\texttt{laws.pd\_snd\_max}).

\subsection{Theorem (certified integration)}
\emph{Domination.} Let $\mathcal H_H$ and $\mathcal H_D$ be the corrupted history and dwell innovations. The $k_H$-th smallest of the history powers is a nondecreasing function of each power, so over all values of the corrupted entries it is smallest when they are zero; hence $s_{k_H}^2\ge\underline s^2$, the value computed with those entries set to zero. Every uncorrupted dwell term is a clean power divided by $s_{k_H}^2\ge\underline s^2$, and every clipped term is at most $c$. Therefore $T\le W$ for all interference values. For binary integration, an uncorrupted exceedance of the actual statistic is an exceedance at $\underline s^2$, and a corrupted term counts at most one. Adding elements to $\mathcal H$ can only lower $\underline s^2$ and replace terms by their maximum, so $W$ is nondecreasing in $\mathcal H$.
\emph{Calibration.} Let the test hit set $\mathcal H$ be independent of the clutter and let $\tilde{\mathcal H}\sim Q$ with $\mathcal H\subseteq\tilde{\mathcal H}$ (a coupling that exists under stochastic dominance). Then $T\le W(X,\mathcal H)\le W(X,\tilde{\mathcal H})$. The pairs $(X_i,\mathcal H_i)$ and $(X,\tilde{\mathcal H})$ are exchangeable, so the split-conformal rank argument gives $\PP\{W(X,\tilde{\mathcal H})>\widehat q\}\le\alpha$. No property of the interference values is used, so the bound holds for adversarial amplitudes chosen after seeing the clutter, as long as the hit positions are independent of the clutter.
\emph{Trimmed sums.} If a corrupted dwell term can take any value, the trimmed sum over the $K-t$ smallest terms is unbounded once more than $t$ dwell terms are corrupted, which a Bernoulli hit model allows with positive probability. The conformal quantile of the worst case is then infinite whenever that probability exceeds $\alpha$.

\subsection{Conformal calibration size and loss}
For continuous i.i.d. scores the $k$-th order statistic of $n$ calibration scores, mapped through the score distribution function, is Beta$(k,n+1-k)$; the conditional false-alarm probability is one minus it, hence Beta$(n+1-k,k)$. The smallest $n$ with $\PP\{P_{\rm fa}>2\alpha\}\le0.05$ is found by search (\texttt{laws.n\_required}). For $P_{\rm d}=P_{\rm fa}^{\gamma}$, $\mathbb EU^\gamma=B(n+1-k+\gamma,k)/B(n+1-k,k)$.
